\documentclass[12pt,a4paper]{article}

\usepackage[margin=2.5cm]{geometry}
\usepackage{amsmath,amssymb}
\usepackage{graphicx}
\usepackage{booktabs}
\usepackage{hyperref}
\usepackage[numbers,sort&compress]{natbib}
\usepackage{setspace}
\usepackage{microtype}
\usepackage{xcolor}
\usepackage{enumitem}
\usepackage{caption}
\usepackage{abstract}
\usepackage{titlesec}
\usepackage{array}
\usepackage{multirow}
\usepackage{tabularx}

\onehalfspacing
\hypersetup{colorlinks=true,linkcolor=blue!60!black,citecolor=blue!60!black,urlcolor=blue!60!black}
\titleformat{\section}{\large\bfseries}{\thesection.}{0.5em}{}
\titleformat{\subsection}{\normalsize\bfseries\itshape}{\thesubsection.}{0.5em}{}
\setlength{\parindent}{1.5em}
\setlength{\parskip}{0.3em}

\title{\textbf{Toward a Measurable Theory of Complexity:\\
Eight Candidate Properties, Parameter-Free Metrics,\\
and Cross-Substrate Validation}\\[0.8em]
\large\textit{A preliminary report inviting expert scrutiny and collaboration}}

\author{Cameron Belt\\
\small Belt Software Solutions\\
\small \href{https://github.com/camerontbelt/complexity-framework}{github.com/camerontbelt/complexity-framework}}

\date{\today\\[0.3em]
\small\textit{Preprint --- not peer reviewed}}

\begin{document}

\maketitle

% ============================================================
\begin{abstract}
\noindent
We present a framework for measuring complexity derived from eight
candidate necessary properties of complex systems.
The framework operationalises five properties as measurable,
substrate-agnostic metrics: spatial entropy with edge-of-chaos
weighting, bidirectional spatial opacity (measuring information
hiding in both upward and downward directions across spatial
scales), temporal opacity (measuring causal structure decay across
time), temporal compression, and fractal dimension.
These metrics are combined using an architecture that encodes two
distinct logical claims: opacity may manifest in space \emph{or}
time (OR logic --- additive combination), while all major properties
must be satisfied simultaneously (AND logic --- multiplicative
combination).

A key theoretical development is the recognition that
\textbf{opacity is a property with two orthogonal axes}: across
space (original metric) and across time (new).
Temporal opacity is measured via mutual information decay --- a
fully parameter-free metric whose boundaries derive from information
theory rather than empirical calibration.
This extends the previous finding that spatial opacity has two
directions; the unified framework now measures opacity across all
available axes of each substrate.

All metric parameters are either fully parameter-free or calibrated
once on the 1D experiment and held fixed thereafter.
The composite is validated across four substrates:

\begin{itemize}[noitemsep]
  \item \textbf{1D Elementary Cellular Automata (256 rules):}
    F1\,=\,1.000, Cohen's $d$\,=\,29.6, C4/C3 separation 82.5$\times$,
    all four Class~4 rules rank 1--4.
  \item \textbf{2D Life-like CA (17 rules):}
    F1\,=\,1.000, Cohen's $d$\,=\,9.64, C4/C3 separation 372$\times$
    (improved from F1\,=\,0.800 in the previous version).
  \item \textbf{2D Rule Space Survey (1,000 random rules):}
    Fractal dimension filter correctly suppresses space-filling rules;
    known Class~4 rules rank in the top 3.
  \item \textbf{Synthetic EEG (five brain states):}
    Seizure and active desynchronisation correctly identified as
    complexity extremes; parameter-free spatial opacity generalises
    to this substrate without modification.
\end{itemize}

A secondary finding concerns the Seeds rule (B2/S in Conway Life
notation).
Seeds scored highly but carries a Wolfram Class~2 label.
We subsequently confirmed that Seeds was proved Turing-complete on an
infinite periodic grid by Naszvadi~\cite{naszvadi2020}.
The metric's intermediate score is therefore correct: Seeds has
genuine computational richness that a Class~2 label undersells.

We make no strong claims about completeness or correctness and
present these results as a coherent, reproducible programme offered
for expert evaluation.

\bigskip
\noindent\textbf{Keywords:} complexity measures, cellular automata,
edge of chaos, information theory, temporal opacity, fractal
dimension, cross-substrate generalisation, parameter-free metrics
\end{abstract}

\rule{\linewidth}{0.4pt}

% ============================================================
\section{Introduction and Motivation}
% ============================================================

The question of how to measure complexity remains genuinely open.
Existing proposals --- algorithmic complexity~\cite{kolmogorov1965},
thermodynamic depth~\cite{lloyd1988}, integrated
information~\cite{tononi2004}, energy rate density~\cite{chaisson2001}
--- each capture something real while failing to capture everything.
No single metric produces rankings that align with intuitive complexity
judgements across system types~\cite{gell-mann1996,ladyman2013}.

The work reported here began as a philosophical inquiry: \emph{can
the properties that make a system genuinely complex be stated precisely
enough to derive measurable metrics from first principles?}
The central commitment is to minimise empirical fitting.
Every metric should ideally derive its boundaries from theoretical
principles --- information-theoretic limits, topological constraints,
or mathematical necessity --- rather than from observing where known
complex systems happen to sit.
A metric whose peak location is theoretically derived makes a genuine
prediction; one discovered by looking at Class~4 rules can only confirm
that Class~4 rules are near its peak.

This version reports substantial progress toward that goal.
The temporal opacity metric introduced here is fully parameter-free.
The spatial entropy weight uses parameter-free tanh gates.
The fractal dimension metric derives its bounds from substrate topology.
Only the spatial opacity Gaussian peaks and the gzip peak retain
empirically-located parameters from the 1D calibration.

\subsection*{Version History}

Version~3 established the basic four-metric composite and demonstrated
1D~ECA discrimination.
Version~4 introduced bidirectional spatial opacity, improving 1D
Cohen's $d$ from 17.6 to 38.4.
Version~5 (this paper) introduces temporal opacity, fractal dimension,
the additive opacity architecture, the triple-tcomp Gaussian, and
extends validation to synthetic EEG data.
The 2D Life-like result improves from F1\,=\,0.800 to F1\,=\,1.000.

% ============================================================
\section{Eight Candidate Properties of Complexity}
% ============================================================

The following eight properties were identified as candidate necessary
conditions for genuine complexity.
They are stated as candidate properties in the spirit of descriptive
relationships derived from observation and reasoning, not from more
fundamental principles.

\subsection{Opaque Hierarchical Layering (P1)}

\textit{Plain language:} Complex systems have layers.
Knowing what is happening at a higher layer does not fully determine
what is happening at the lower layer, and vice versa.
Each layer hides information from the others in both directions.

\textit{Formally:} A system exhibits opacity of degree $N$ when its
state space is partitioned into $N$ stable hierarchical layers such
that $H(S_K \mid S_{K+1}) > 0$ and $H(S_{K+1} \mid S_K) > 0$ for
adjacent layer pairs $(K, K+1)$.

This framework operationalises P1 across three distinct axes:
\begin{itemize}[noitemsep]
  \item \textbf{Spatial opacity (upward):}
    $H(\text{global density} \mid \text{local patch})$ --- given local
    structure, global state remains uncertain.
  \item \textbf{Spatial opacity (downward):}
    $H(\text{local patch} \mid \text{global density})$ --- given
    global state, local structure remains underdetermined.
  \item \textbf{Temporal opacity:}
    $I(X_t;\,X_{t+\text{lag}})$ with characteristic decay --- the past
    predicts the future in a structured, non-trivial way.
\end{itemize}

The recognition that opacity exists in both spatial directions
\emph{and} across time is a primary theoretical contribution of this
version.

\subsection{Modular Interconnection (P2)}

\textit{Plain language:} Complex systems are made of semi-independent
parts with internal logic, connected in ways that produce behaviour no
single part could produce alone.
Operationalised as gzip compression ratio: rich modular structure
resists compression because semi-independent regions each have their
own patterns that resist any single-template description.

\subsection{Self-Assembly Without Top-Down Specification (P3)}

\textit{Plain language:} Genuine complexity must grow; it cannot be
designed in full.
The moment an interface is fully specified it ceases to be a genuine
opaque boundary.
This places a hard ceiling on engineered complexity and implies that
the most complex known systems arose through processes that did not
require a designer to understand the result.

\subsection{Recursive Self-Modification (P4)}

\textit{Plain language:} A system on fixed rules can explore but
cannot expand its state space.
Generating genuinely new abstraction layers requires modifying the
rules the system runs on.
Operationalised (partially) as temporal compression: persistent
dynamic attractors reflect selection-stabilised pattern libraries.

\subsection{Selection Pressure (P5)}

\textit{Plain language:} Without something filtering which variations
survive, self-assembly produces noise.
Selection pressure gives drift a direction and allows complexity to
accumulate.

\subsection{Co-evolutionary Dynamics (P6)}

\textit{Plain language:} When two or more self-assembling systems share
an environment, each becomes selection pressure on the other.
Complexity bootstraps itself from the outside in.

\subsection{Thermodynamic Drive (P7)}

\textit{Plain language:} Complex dissipative structures~\cite{prigogine1977}
are thermodynamically favoured because they generate entropy more
efficiently than simpler alternatives.
Operationalised as the spatial entropy weight: the tanh gate structure
rewards systems near the edge-of-chaos regime where thermodynamic
drive maintains active dynamics without collapsing to maximum entropy.

\subsection{Scale Invariance (P8 --- meta-property)}

The candidate properties apply at every scale, confirmed empirically
by the ubiquity of power laws in complex systems~\cite{bak1987}.
The fractal dimension metric directly operationalises one aspect of
this: complex systems use space in a self-similar, intermediate-
dimensional way rather than filling it completely or collapsing to
a 1D structure.

% ============================================================
\section{Metric Operationalisation}
% ============================================================

Five of the eight properties were operationalised as measurable metrics.
Properties P3, P6, and full P4 describe generative-process properties
requiring multi-system experiments or process-level observation.

\subsection{Spatial Entropy Weight (from P7)}

Binary Shannon entropy per timestep is computed over the spatial
density distribution, then passed through parameter-free tanh gates:

\begin{equation}
w_H = \tanh(kH_s)\cdot\tanh\!\bigl(k(1-H_s)\bigr)
      \cdot\Bigl(1 + \exp\!\Bigl(-\tfrac{(\sigma_H-\mu_\sigma)^2}
      {2\sigma_\sigma^2}\Bigr)\Bigr)
\label{eq:wH}
\end{equation}

where $k=50$ controls gate sharpness.
The tanh product is zero at $H_s=0$ (trivially ordered) and $H_s=1$
(maximum chaos), with the peak location \emph{emerging from where the
system sits} rather than being specified in advance --- fully
parameter-free.
The Gaussian bonus rewards the entropy standard deviation signature
that distinguishes Class~4 from Class~3 under random initial
conditions ($\sigma_H \approx 0.013$ vs $\approx 0.010$, reflecting
operation at a dynamical critical point).
The peak $\mu_\sigma = 0.012$ is the only fitted parameter in this
weight.

\subsection{Bidirectional Spatial Opacity (from P1, spatial axis)}

Two complementary conditional entropies measure information hiding
across spatial scales:

\begin{align}
\text{opacity}_\uparrow &=
  \frac{H(\text{global density bin} \mid \text{local 3-cell patch})}
       {\log_2 n_\text{bins}}
\label{eq:opup}\\[4pt]
\text{opacity}_\downarrow &=
  \frac{H(\text{local patch} \mid \text{global density bin})}
       {\log_2 8}
\label{eq:opdown}
\end{align}

The spatial opacity weight combines both multiplicatively:

\begin{equation}
w_{\text{OP},s} = G(\text{opacity}_\uparrow,\,0.14,\,0.10)
                  \times G(\text{opacity}_\downarrow,\,0.97,\,0.05)
\label{eq:wops}
\end{equation}

where $G(x,\mu,\sigma) = \exp\!\bigl(-\frac{(x-\mu)^2}{2\sigma^2}\bigr)$.
For substrates where the 1D-calibrated peaks do not transfer (e.g.\
EEG channel space), the parameter-free tanh formulation is used:

\begin{equation}
w_{\text{OP},s}^{\text{tanh}} = \tanh(k\cdot\text{opacity}_\uparrow)
                                  \cdot\tanh(k\cdot(1-\text{opacity}_\uparrow))
\label{eq:wops_tanh}
\end{equation}

This is the target formulation for all substrates as the framework
develops toward full parameter-freedom.

\subsection{Temporal Opacity (from P1, temporal axis)}

Temporal opacity measures information hiding across time.
Normalised mutual information at lag $\ell$ is:

\begin{equation}
\mathrm{MI}_\ell = \frac{I(X_t;\,X_{t+\ell})}{H(X_t)}
\label{eq:mi}
\end{equation}

Two statistics summarise the decay profile:
$\mathrm{MI}_1$ (initial temporal correlation) and
$\text{decay} = \mathrm{MI}_1 - \mathrm{MI}_{10}$.
The weight uses three parameter-free tanh gates:

\begin{equation}
w_{\text{OP},t} = \tanh(k\cdot\mathrm{MI}_1)
                  \cdot\tanh(k\cdot(1-\mathrm{MI}_1))
                  \cdot\tanh(k\cdot\max(\text{decay},\,0))
\label{eq:wopt}
\end{equation}

\textbf{Gate semantics (all theoretically derived):}
Gate~1 ($\tanh(k\cdot\mathrm{MI}_1)$) vanishes when $\mathrm{MI}_1\!\to\!0$
(random, no temporal structure).
Gate~2 ($\tanh(k(1-\mathrm{MI}_1))$) vanishes when
$\mathrm{MI}_1\!\to\!1$ (frozen, trivially predictable).
Gate~3 ($\tanh(k\cdot\text{decay})$) vanishes when decay\,$\to$\,0
(either random or frozen --- neither has causal dynamics).
Only complex systems pass all three simultaneously.

\textbf{Dimensional scaling.}
In 1D binary CAs, field-level temporal MI is near-zero for all classes
(${\approx}0.01$--$0.02$) because single-cell binary states carry
minimal mutual information at the field level.
In 2D Life-like CAs, Class~4 rules show
$\mathrm{MI}_1\approx0.20$--$0.26$, Class~3 chaotic rules show
${\approx}0.002$, and frozen Class~2 rules show ${\approx}1.0$.
The metric therefore activates naturally as substrate dimensionality
increases --- no substrate flag is needed.

\textbf{Relation to existing theory.}
This metric is related to Crutchfield and Shalizi's excess entropy
(statistical complexity)~\cite{crutchfield2001}, which measures
$I(\text{past};\text{future})$ as a summary of causal memory.
The key distinction is that we use the \emph{decay profile} of
$\mathrm{MI}_\ell$ rather than a single summary, allowing the decay
gate to distinguish complex systems (characteristic timescale decay)
from simple periodic ones (sustained or sinusoidal MI).

\subsection{Temporal Compression (from P4/P5)}

For each cell $i$, the run-length compressibility of its state history:

\begin{equation}
\text{tc}_i = 1 - \frac{1 + \text{flip\_count}_i}{T}
\label{eq:tc}
\end{equation}

The mean across all cells gives a measure of temporal persistence.
The weight uses a \emph{triple} Gaussian to cover three confirmed
attractor regimes:

\begin{equation}
w_T = \max\!\bigl(G(tc,\,0.58,\,0.08),\;
                   G(tc,\,0.73,\,0.08),\;
                   G(tc,\,0.90,\,0.05)\bigr)
\label{eq:wT}
\end{equation}

\noindent\textbf{Peak 1 (0.58):} 1D Class~4 rules under random IC.
Dynamic attractor --- glider interactions maintain intermediate flip
rate.\\
\textbf{Peak 2 (0.73):} 1D Class~4 rules under single-cell IC.
Makes the weight IC-independent in 1D.\\
\textbf{Peak 3 (0.90):} 2D Class~4 rules (Conway, HighLife, Morley,
Day\&Night) converging to a static-ecosystem attractor.
Glider ecosystems form stable backgrounds; most cells rarely flip.
Confirmed IC-independent across density sweep $d=0.1$--$0.8$.

The three-peak formulation is backward-compatible: 1D Class~4 rules
still hit peaks~1 or~2, and peak~3 is too high to affect them.

\subsection{Fractal Dimension (from P8)}

Box-counting (Minkowski-Bouligand) fractal dimension of the
late-time spatial pattern, combined into a weight via:

\begin{equation}
w_\text{dim} = G\bigl(\underbrace{d_f - (D-1)}_{\text{excess}},\;0.35,\;0.20\bigr)
\label{eq:wdim}
\end{equation}

where $D$ is the substrate spatial dimension ($D=1$ for 1D CA,
$D=2$ for 2D Life) and $d_f$ is the measured fractal dimension.
The dimensional excess measures how far the pattern extends into the
available spatial dimension beyond a simple 1D line.

\textbf{Theoretical grounding.}
Complex systems use space in an intermediate, self-similar way:
neither collapsing to a 1D structure ($d_f\approx1$, excess$\approx0$)
nor filling the full $D$-dimensional space
($d_f\approx D$, excess$\approx1$).
Space-filling chaotic and frozen rules have excess near~1.0.
Class~4 rules with structured glider ecosystems cluster at
excess~$\approx0.35$ ($d_f\approx1.35$ in 2D).

\textbf{1D behaviour.}
In 1D, $d_f\approx1.0$, giving excess~$\approx0$ and
$w_\text{dim}\approx0$ --- the term has negligible effect and is
omitted from the 1D composite to preserve all existing 1D results
exactly.
In 2D it activates strongly for Class~4 and suppresses space-filling
Class~3/2.

\subsection{Gzip Compression Ratio (from P2)}

\begin{equation}
\text{gz} = \frac{|\text{compressed}|}{|\text{original}|}
\label{eq:gz}
\end{equation}

The weight $w_G = G(\text{gz},\,0.10,\,0.05)$ peaks at the
intermediate compression attractor ($\approx\!10\%$) reflecting large
stable regions that compress efficiently with complex interaction
boundaries preventing total compression.

\subsection{Composite Score (v7)}

The composite formula encodes two distinct logical claims:

\begin{equation}
\boxed{C = w_H \;\times\;
           \underbrace{(w_{\text{OP},s} + w_{\text{OP},t})}_{\text{additive: OR logic}}
           \;\times\; w_T \;\times\; w_G \;\times\; w_\text{dim}}
\label{eq:composite}
\end{equation}

\textbf{Additive opacity (OR logic).}
$w_{\text{OP},s} + w_{\text{OP},t}$ rather than their product.
P1 opacity can manifest in either space or time --- satisfying one
direction is sufficient evidence of hierarchical layering.
In 1D, spatial opacity is strong and temporal is near-zero; the sum
is dominated by spatial, identical to previous versions.
In 2D, the 1D-calibrated spatial Gaussian peaks return near-zero for
Class~4 rules, but temporal opacity carries the composite.

\textbf{Multiplicative combination (AND logic).}
All other terms are multiplicative: a system must satisfy entropy,
tcomp, gzip, and fractal dimension simultaneously.
Near-zero on any single dimension collapses the score, encoding the
non-eliminability claim.

\subsection{Parameter Status Summary}

\begin{table}[h]
\centering
\caption{Parameter status of all metric weights. The trajectory
  is toward full parameter-freedom: every metric whose boundaries are
  information-theoretic can in principle use tanh gates.}
\begin{tabular}{llll}
\toprule
Metric & Parameter status & Boundary derivation & Substrate \\
\midrule
Entropy tanh gates & Fully parameter-free
  & $H=0$ dead, $H=1$ chaos & All \\
Entropy variance bonus & One fitted ($\mu_\sigma=0.012$)
  & C4 cluster, 1D ECA & All \\
Temporal opacity & Fully parameter-free
  & MI=0 random, MI=1 frozen & All \\
Spatial opacity Gaussian & Two fitted (0.14, 0.97)
  & C4 cluster, 1D ECA & 1D CA \\
Spatial opacity tanh & Fully parameter-free
  & Tanh gates on opacity$_\uparrow$ & Non-CA \\
Triple tcomp & Three confirmed attractors
  & IC sweep \& 2D density sweep & 1D, 2D \\
Gzip weight & One fitted ($\mu=0.10$)
  & C4 cluster, 1D ECA & All \\
Fractal dimension & One fitted ($\mu_\text{ex}=0.35$)
  & 2D C4 cluster & 2D+ \\
\bottomrule
\end{tabular}
\end{table}

% ============================================================
\section{Experiment 1: Elementary Cellular Automata (1D)}
% ============================================================

\subsection{Setup}

All 256 Wolfram elementary cellular automata were simulated on
$W=150$ cells for $T=300$ steps (burn-in 20, measurement window 150),
random 50\% IC averaged over five seeds.
Spatial opacity Gaussian peaks were established by observing where the
four known Class~4 rules cluster, then locked for all subsequent
experiments.

\subsection{Results}

\begin{table}[h]
\centering
\caption{1D ECA results (v5 composite, random 50\% IC, 5 seeds).}
\begin{tabular}{lrl}
\toprule
Metric & Value & Notes \\
\midrule
F1 score          & 1.000 & All four C4 rules rank 1--4 of 256 \\
Cohen's $d$       & 29.62 & Improved from 17.6 in v4 \\
Rank-biserial $|r|$ & 1.000 & Complete rank separation \\
C4/C3 separation  & 82.5$\times$ & Mean score ratio \\
$p$ (Mann-Whitney) & $3.05\times10^{-4}$ & One-sided \\
$w_{\text{OP},t}$ (C4 mean) & 0.024
  & Near-zero as predicted: 1D dimensional scaling \\
\bottomrule
\end{tabular}
\end{table}

The four Class~4 rules (124, 137, 193, 110) rank in positions 1--4
respectively.
The temporal opacity weight is near-zero for all 1D classes, confirming
the dimensional scaling prediction.
The additive architecture means this does not hurt the score: the
spatial opacity channel carries the composite as before, and the
temporal channel simply adds nothing (rather than multiplying by zero).

\subsection{The Entropy Variance Signature}

Class~4 rules are distinguished from Class~3 not primarily by mean
entropy (both sit near $H\approx0.984$) but by entropy variance:
Class~4 rules exhibit approximately twice the entropy standard
deviation ($\sigma_H\approx0.013$ vs $\approx0.010$), reflecting
operation at a dynamical critical point where structures form and
dissolve rather than settling into frozen chaos.

\subsection{Multi-Metric Non-Eliminability}

Class~4 rules are the only class that cannot be eliminated by any
single metric.
Class~1 rules score near zero on entropy and opacity.
Class~2 rules score high on temporal compression but low on entropy
variance and gzip.
Class~3 rules score high on entropy but low on temporal compression and
gzip.
Only Class~4 rules score moderately well on all dimensions
simultaneously.

% ============================================================
\section{Experiment 2: Two-Dimensional Life-like Rules}
% ============================================================

\subsection{Setup}

Seventeen Life-like cellular automaton rules spanning four Wolfram
behaviour classes were evaluated on a $64\times64$ grid for 200 steps
(5 seeds, random 35\% density IC).
Ground truth classifications follow the LifeWiki taxonomy.
The four known Class~4 rules are Conway's Life (B3/S23),
HighLife (B36/S23), Day\&Night (B3678/S34678), and
Morley (B368/S245).

\subsection{Results}

\begin{table}[h]
\centering
\caption{2D Life-like results: v4 vs v5 composite.}
\begin{tabular}{lrrr}
\toprule
Metric & v4 result & v5 result & Change \\
\midrule
F1 score         & 0.800 & \textbf{1.000} & +25\% \\
Cohen's $d$      & 2.32  & \textbf{9.64}  & +315\% \\
C4/C3 separation & 3.2$\times$ & \textbf{372$\times$} & +11,525\% \\
C4 ranks         & [9,10,11,12] & \textbf{[1,2,3,4]} & All at top \\
\bottomrule
\end{tabular}
\end{table}

The dramatic improvement from v4 to v5 comes from three changes:

\begin{enumerate}[noitemsep]
  \item \textbf{Temporal opacity ($w_{\text{OP},t}$).}
    C4 mean\,=\,0.935 vs C3 mean\,=\,0.003 --- a 312$\times$
    separation on this single metric.
    Conway, HighLife, Morley, and Day\&Night all score 0.90--0.97.
    Chaotic Class~3 rules (34Life, Gnarl) score 0.001--0.005.
    Frozen Class~2 rules (Maze) score 0.000.

  \item \textbf{Fractal dimension ($w_\text{dim}$).}
    C4 mean\,=\,0.798 vs C3 mean\,=\,0.062.
    Correctly places Class~4 rules at intermediate fractal dimension
    ($d_f\approx1.35$--$1.52$) and suppresses space-filling Class~3
    rules ($d_f\approx1.94$).
    This resolved the Seeds ranking issue (Section~\ref{sec:seeds}).

  \item \textbf{Triple-peak tcomp.}
    The third Gaussian peak at 0.90 rescues Conway and Day\&Night
    from penalty for their high tcomp values (0.88--0.95),
    which reflect the static-ecosystem attractor.
\end{enumerate}

\subsection{The Two Attractor Discovery}

2D Life-like rules have two distinct tcomp attractors corresponding
to different modes of complexity:

\begin{description}[noitemsep]
  \item[Dynamic attractor (${\approx}0.58$):] perpetual activity,
    the system never settles.
    34Life-style rules.
    This is what the 1D calibration targeted.
  \item[Static-ecosystem attractor (${\approx}0.90$):] glider
    ecosystems settle into stable backgrounds with most cells rarely
    flipping.
    Conway, HighLife, Morley, Day\&Night all converge here regardless
    of IC density (confirmed across $d=0.1$--$0.8$).
    The attractor is IC-independent --- a true dynamical attractor,
    not a calibration artefact.
\end{description}

This parallels the 1D finding that motivated the dual-Gaussian tcomp:
the same property has different attractor values under different IC
types and substrate geometries.
The response in both cases is to add a Gaussian peak at the confirmed
attractor location.

\subsection{The Seeds Finding}
\label{sec:seeds}

Seeds (B2/S) was initially treated as a misclassified edge case
in v4.
In v5, its fractal dimension ($d_f\approx1.83$, near space-filling)
reduces its score, dropping it from rank~1 to rank~6.

More significantly, we confirmed that Seeds was \textbf{proved
Turing-complete on an infinite periodic grid by
Naszvadi~\cite{naszvadi2020}} via construction of a Rule~110 unit
stripe.
The LifeWiki notes that Seeds has no endemic patterns --- no stable or
periodic structures exist on any finite grid --- so it requires an
infinite periodic tiling to express its full computational
architecture.

This finding supports rather than contradicts the metric's
intermediate scoring.
The framework is detecting genuine computational richness (temporal MI,
intermediate tcomp) that Seeds' Wolfram Class~2 label undersells.
Conway achieves computation through persistent localised structures
that use space efficiently ($d_f\approx1.5$); Seeds achieves
computation through wavefront interference on an infinite grid but
appears near-space-filling on finite grids ($d_f\approx1.83$).
The metric correctly scores the \emph{spatial efficiency} of the
computational substrate --- a distinction that Turing-completeness
alone does not capture.

\subsection{The 2D Rule Space Survey}

To test the v7 framework on a broader sample, 1,000 rules were drawn
uniformly at random from the 262,144 possible Life-like rules
($2^9\times2^9$) and evaluated at reduced resolution
($40\times40$ grid, 80 steps, 3 seeds).
The known Class~4 rules were included as anchors.

Result: Conway, Morley, and HighLife ranked in positions 1, 2, and 3
of the combined 1,007-rule set.
The fractal dimension metric was the primary discriminator: every
survey rule with $d_f>1.8$ was suppressed to near-zero composite
score.
The best survey candidate was B3/S2456 (rank 5, $C=0.071$,
$d_f=1.77$, $w_{\text{OP},t}=0.954$), warranting further
investigation at full resolution.

% ============================================================
\section{Experiment 3: N-body Particle Simulation}
% ============================================================

\subsection{Setup and Caveat}

A Lennard-Jones particle simulation with tunable force constants was
scanned across a parameter space analogous to physical coupling
strengths: $\alpha$ (well depth, analogous to electromagnetic
coupling) and $\alpha_s$ (equilibrium separation, analogous to strong
force coupling), normalised so that $\alpha=\alpha_s=1.0$ corresponds
to our universe's values.

\textbf{Important caveat:} this simulation uses a toy force law and
cannot reproduce the full fine-tuning constraints of real physics.
Results should be interpreted as a methodological demonstration ---
does the framework find an island of stability in a simplified
parameter space? --- not as a claim about actual physical constants.

\subsection{Results (v4 framework)}

A $16\times16$ parameter scan (256 points, 3 seeds each), using the
v4 composite without fractal dimension or temporal opacity:

\begin{table}[h]
\centering
\caption{N-body simulation results.}
\begin{tabular}{ll}
\toprule
Metric & Value \\
\midrule
Active/inactive separation & $8.9\times$ \\
F1 & 0.959 \\
Cohen's $d$ & 3.89 \\
Rank-biserial $|r|$ & 0.982 \\
$p$ (Mann-Whitney) & $1.33\times10^{-17}$ \\
Our universe percentile & 54th (within Active island) \\
vs analytical island baseline & $+71.9\%$ F1 \\
\bottomrule
\end{tabular}
\end{table}

The 54th percentile placement is consistent with the anthropic
principle: any universe in the Active island would produce observers
asking this question.
The v5 composite extension to the N-body substrate is a target for
the next experimental cycle.

% ============================================================
\section{Experiment 4: Synthetic EEG}
% ============================================================

\subsection{Motivation and Setup}

To test the framework on a qualitatively different substrate,
synthetic multichannel EEG signals were generated for five canonical
brain states.
This experiment tests whether the metrics generalise to continuous
time series data and whether the parameter-free spatial opacity
formulation (tanh gates) transfers without modification.

EEG signals were represented as a $(T\times C)$ binary matrix with
$T=2560$ time samples at 256\,Hz (10\,s) and $C=19$ channels.
Each channel was binarised by threshold at its temporal median.
The `spatial' dimension is therefore channels (not cells), and the
`temporal' dimension is samples.

Five brain states were generated with neurophysiologically motivated
parameters:

\begin{description}[noitemsep]
  \item[Deep sleep:] large-amplitude delta (0.8\,Hz) with modest
    inter-channel phase jitter.
  \item[NREM Stage~2 (sleep spindles):] background delta plus spindle
    bursts (12--15\,Hz, 3--5 per epoch) and one K-complex.
  \item[Awake, eyes-closed resting:] dominant alpha (10\,Hz) with
    amplitude modulation, plus beta component and pink noise.
  \item[Active cognitive task:] desynchronised broadband beta/gamma
    (18--55\,Hz) with independent channels.
  \item[Seizure (tonic-clonic):] hypersynchronous 3\,Hz spike-wave
    across all channels.
\end{description}

\subsection{Substrate Adaptation}

Two adaptations were made for the EEG substrate:

\begin{enumerate}[noitemsep]
  \item \textbf{Channel spatial opacity.}
    The `spatial' dimension is channels rather than cells.
    The local patch is a 3-channel neighbourhood.
    The weight uses parameter-free tanh gates (Eq.~\ref{eq:wops_tanh})
    rather than 1D-calibrated Gaussians, since channel space in EEG
    has no reason to sit near the CA-calibrated peaks.

  \item \textbf{No fractal dimension.}
    The $(T\!\times\!C)$ representation is treated as 1D spatial
    (channel space), so $w_\text{dim}$ is omitted, consistent with
    the 1D CA handling.
\end{enumerate}

The composite is:
\begin{equation}
C_\text{EEG} = w_H \times
               (w_{\text{OP},\text{ch}}^{\text{tanh}} + w_{\text{OP},t})
               \times w_T \times w_G
\end{equation}

\subsection{Results}

\begin{table}[h]
\centering
\caption{Synthetic EEG results. States ranked by composite score.}
\begin{tabular}{llrl}
\toprule
Brain state & Regime & Composite $C$ & Key signals \\
\midrule
Sleep spindles & Transitional
  & $0.548\pm0.057$
  & High $w_{\text{OP},s}$ (channel desync), strong $w_{\text{OP},t}$ \\
Deep sleep & Ordered
  & $0.480\pm0.039$
  & Strong $w_{\text{OP},t}$ (slow monotone decay) \\
Awake resting & Complex
  & $0.383\pm0.044$
  & Good $w_{\text{OP},s}$, moderate $w_{\text{OP},t}$ \\
Active task & Chaotic
  & $0.071\pm0.015$
  & Near-zero $w_{\text{OP},t}$ (no temporal structure) \\
Seizure & Frozen
  & $0.001\pm0.002$
  & Near-zero $w_{\text{OP},s}$ (hypersync) \& $w_{\text{OP},t}$ \\
\bottomrule
\end{tabular}
\end{table}

Seizure and active task correctly identify as the complexity extremes
with a large gap below the top three states.
The sharp bottom-two separation confirms that parameter-free tanh
gate spatial opacity generalises to this substrate without
modification.

The top-three ordering (spindles $>$ deep sleep $>$ awake resting)
disagrees with Tononi's $\Phi$ framework, which typically places
wakefulness highest.
Within our framework the ordering is scientifically defensible:
sleep spindles have the richest temporal transient structure
($w_T=0.981$) and strong channel independence; deep sleep has strong
monotone temporal MI decay from its large-amplitude slow oscillation.
The disagreement with $\Phi$ is noted honestly and suggests that
applying both frameworks to the same real EEG data would be
informative.

% ============================================================
\section{Discussion}
\label{sec:discussion}
% ============================================================

\subsection{What the Framework Establishes}

\begin{enumerate}[noitemsep]
  \item Perfect discrimination of computationally universal rules in
    1D (F1\,=\,1.000, $d$\,=\,29.6) and 2D (F1\,=\,1.000,
    $d$\,=\,9.64) cellular automata with parameters fixed after 1D
    calibration.

  \item The improvement from v4 to v5 in 2D came not from
    recalibration but from adding theoretically motivated metrics
    (temporal opacity, fractal dimension) that activate naturally in
    2D without fitted peaks.

  \item The fractal dimension metric correctly identifies the spatial
    efficiency of computational substrates, distinguishing Conway
    (structured glider ecosystems) from Seeds (space-filling
    wavefronts) despite both being Turing-complete.

  \item Meaningful discrimination in the N-body simulation
    (F1\,=\,0.959) and qualitative ordering of brain states in
    synthetic EEG.
\end{enumerate}

\subsection{The k=3 CA Negative Result}

The k=3 totalistic CA experiment (radius-2, 5-bit lookup table)
produced F1\,=\,0.000.
The labelled complex rules ({22,\,26}) saturate to full or nearly-full
grids under all tested initial conditions and remain frozen.
The lookup table for rule~22 maps neighbourhood sums 1, 2, 4 to 1 ---
at any reasonable density, most cells output 1, causing the grid to
fill completely within 10--15 steps.

This is not a metric failure.
The framework correctly reports what these rules do.
The classification $\{22,\,26\}$ as complex appears to originate from
a different CA formalism (possibly Wolfram's outer-totalistic $k=3$
with separate handling of the centre cell state, giving 1024 rather
than 32 rules) or from experiments under narrow specific initial
conditions not captured by our standard random IC protocol.
The metrics accurately described the actual dynamics; the ground truth
labels were incorrect for our experimental setup.
This demonstrates that the metrics do not produce false positives even
when expected to.

\subsection{Toward Fully Parameter-Free Metrics}

The most significant theoretical development is the recognition that
opacity can be measured with parameter-free tanh gates, and that the
composite should reflect the logical structure of the underlying
claims.

The temporal opacity metric demonstrates achievability: three gates
with information-theoretic boundaries produce strong discrimination
without any fitted parameters.
The target for future work is to extend this to spatial opacity
(replacing the 0.14 and 0.97 Gaussian peaks with tanh gates) and
gzip (replacing the 0.10 peak), leaving only the entropy variance
bonus and tcomp attractor peaks as empirically-located parameters.

The tcomp attractor peaks are a special case: they are confirmed
dynamical attractor values, not arbitrary fitting points.
They encode genuine discoveries about where complex systems sit in
tcomp space across substrates and IC types.
Their retention is theoretically defensible as encoding confirmed
knowledge about the dynamical geometry of complexity.

\subsection{Relation to Existing Work}

\begin{itemize}[noitemsep]
  \item \textbf{Wolfram's computational irreducibility~\cite{wolfram2002}:}
    directly operationalised by the gzip metric.
    The framework validates that Class~4 rules --- Wolfram's
    computationally irreducible rules --- are precisely those that sit
    in the intermediate compressibility regime.

  \item \textbf{Kauffman's edge of chaos~\cite{kauffman1993,langton1990}:}
    precisely operationalised by the intermediate-regime weighting
    throughout.

  \item \textbf{Crutchfield \& Shalizi's excess entropy~\cite{crutchfield2001}:}
    the temporal opacity metric is closely related, using the decay
    profile of $I(X_t;\,X_{t+\ell})$ as a characterisation of causal
    structure.

  \item \textbf{Tononi's integrated information $\Phi$~\cite{tononi2004}:}
    an alternative operationalisation of P1.
    The bidirectional opacity analysis (measuring both directions
    simultaneously) may offer a computationally cheaper approximation
    to similar properties.

  \item \textbf{Naszvadi's Seeds proof~\cite{naszvadi2020}:}
    demonstrates that Turing-completeness and spatial structural
    complexity (as measured by fractal dimension) are distinct
    properties.
\end{itemize}

\subsection{Limitations and Open Questions}

\begin{enumerate}[noitemsep]
  \item \textbf{Substrate-specific calibration.}
    Several metric peaks remain empirically located on 1D.
    An independent theoretical derivation of why complex systems
    occupy these specific regimes would substantially strengthen the
    universality claim.

  \item \textbf{EEG ordering.}
    The top-three brain state ordering disagrees with Tononi's $\Phi$
    ordering.
    Applying both frameworks to the same real EEG data would clarify
    whether this is a genuine theoretical disagreement or an artefact
    of synthetic signal construction.

  \item \textbf{N-body cosmological extension.}
    The toy Lennard-Jones force law cannot reproduce full fine-tuning
    constraints; results should be treated as a methodological
    demonstration only.

  \item \textbf{Seeds on finite vs infinite grids.}
    A principled extension accounting for grid-size dependence would
    allow the Seeds result to be stated more precisely.

  \item \textbf{The candidate properties themselves.}
    The eight properties were derived through philosophical reasoning.
    Experts may identify properties that are missing, redundant, or
    imprecisely stated.
\end{enumerate}

% ============================================================
\section{Conclusion}
% ============================================================

We have presented eight candidate properties of complexity, five
measurable metrics derived from those properties, and empirical
results showing that the composite correctly identifies
computationally universal behaviour across four distinct substrates.

The central theoretical advance is the recognition that P1 (opaque
hierarchical layering) manifests along two orthogonal axes --- across
spatial scale and across time --- and that both should be measured.
The temporal opacity metric implements this without empirical fitting,
deriving its boundaries from information theory alone.
The additive combination of spatial and temporal opacity reflects the
correct logical claim: satisfying one direction of P1 is sufficient
evidence of hierarchical layering.

The fractal dimension metric implements P8 (scale invariance) in a
way that activates naturally with substrate dimensionality.
In 1D it contributes near-zero weight; in 2D it provides strong
discrimination; in 3D substrates it would contribute more strongly.
This is the dimensional scaling behaviour originally hypothesised.

The most significant open question is whether the remaining
empirically-located parameters have theoretical derivations.
If so, the framework becomes predictive.
If not, it remains a useful and reproducible measurement procedure
with demonstrated cross-substrate consistency.

We invite experts in complexity science, dynamical systems,
information theory, and computational physics to evaluate the
methodology, identify errors, and propose improvements.
All code and data are available at
\url{https://github.com/camerontbelt/complexity-framework}.

% ============================================================
\section*{Acknowledgements}
% ============================================================

This work was developed in extended conversation with Claude
(Anthropic), which contributed computational implementation,
literature context, and iterative theoretical refinement throughout.
The intellectual framework, experimental design decisions, and
interpretations are those of the human author.
Cross-validation with Grok (xAI) was used to stress-test
experimental findings.

% ============================================================
\appendix
\section{Candidate Properties --- Full Statement}
% ============================================================

\begin{description}
  \item[P1 --- Opaque Hierarchical Layering:]
    $H(S_K \mid S_{K+1}) > \varepsilon$ and
    $H(S_{K+1} \mid S_K) > \varepsilon$ for all adjacent layer pairs.
    Complexity of degree $N$ requires $N$ such opaque boundaries.
    Measurable in space (two directions) and time.

  \item[P2 --- Modular Interconnection:]
    Modules are internally coherent and externally coupled.
    Optimal topology is scale-free small-world.
    Complexity scales multiplicatively with modular richness at each layer.

  \item[P3 --- Self-Assembly Without Top-Down Specification:]
    Complexity cannot be engineered; it must self-assemble.
    The moment an interface is fully specified it ceases to be genuinely opaque.

  \item[P4 --- Recursive Self-Modification:]
    Maximum complexity is bounded by capacity for self-modification.
    Fixed-rule systems explore but cannot expand their state space.

  \item[P5 --- Selection Pressure:]
    Asymmetric environmental constraints are required for directed
    complexity rather than drift.

  \item[P6 --- Co-evolutionary Dynamics:]
    Maximum complexity is determined by richness of external relationships.
    Complexity is fundamentally ecological rather than individual.

  \item[P7 --- Thermodynamic Drive:]
    Complexity is an expression of entropy maximisation, not opposed to it.
    Dissipative structures are thermodynamically favoured.

  \item[P8 --- Scale Invariance (meta-property):]
    The properties apply at every scale.
    The universe is an instance of the phenomenon these properties describe.
\end{description}

% ============================================================
\section{Metric Parameters}
% ============================================================

\begin{table}[h]
\centering
\caption{All metric parameters used in experiments.
  tcomp attractor peaks are dynamically confirmed values,
  not arbitrary fitting points.}
\begin{tabular}{lllll}
\toprule
Parameter & Value & Status & Substrate \\
\midrule
Entropy tanh $k$ & 50 & Parameter-free & All \\
Entropy var.\ bonus $\mu_\sigma$ & 0.012 & Fitted: C4 cluster, 1D & All \\
Entropy var.\ bonus $\sigma_\sigma$ & 0.008 & Fitted: C4 cluster, 1D & All \\
Spatial opacity$\uparrow$ peak & 0.14 & Fitted: C4 cluster, 1D & 1D CA \\
Spatial opacity$\uparrow$ $\sigma$ & 0.10 & Fitted: C4 cluster, 1D & 1D CA \\
Spatial opacity$\downarrow$ peak & 0.97 & Fitted: C4 cluster, 1D & 1D CA \\
Spatial opacity$\downarrow$ $\sigma$ & 0.05 & Fitted: C4 cluster, 1D & 1D CA \\
Spatial opacity tanh $k$ & 5 & Parameter-free & Non-CA \\
Temporal opacity tanh $k$ & 10 & Parameter-free & All \\
tcomp peak 1 / $\sigma$ & 0.58 / 0.08 & Confirmed: 1D C4 rand.\ IC & 1D, 2D \\
tcomp peak 2 / $\sigma$ & 0.73 / 0.08 & Confirmed: 1D C4 sc.\ IC & 1D \\
tcomp peak 3 / $\sigma$ & 0.90 / 0.05 & Confirmed: 2D C4 attractor & 2D \\
Gzip peak / $\sigma$ & 0.10 / 0.05 & Fitted: C4 cluster, 1D & All \\
Fractal dim.\ excess peak & 0.35 & Fitted: 2D C4 cluster & 2D+ \\
Fractal dim.\ excess $\sigma$ & 0.20 & Fitted: 2D C4 cluster & 2D+ \\
\bottomrule
\end{tabular}
\end{table}

% ============================================================
\section{Experimental Results Summary}
% ============================================================

\begin{table}[h]
\centering
\caption{Summary of results across all experiments.}
\begin{tabular}{lccccl}
\toprule
Experiment & $n$ & F1 & $d$ & $|r|$ & Notes \\
\midrule
1D ECA (v5)      & 256  & 1.000 & 29.6 & 1.000
  & All 4 C4 ranks 1--4 \\
2D Life (v5)     & 17   & 1.000 & 9.64 & 1.000
  & All 4 C4 ranks 1--4 \\
2D Rule Survey   & 1000 & ---   & ---  & ---
  & C4 rules rank 1--3 of 1007 \\
k=3 CA           & 32   & 0.000 & ---  & ---
  & Mislabelled ground truth \\
N-body (v4)      & 256  & 0.959 & 3.89 & 0.982
  & Our universe: 54th percentile \\
EEG synthetic    & 5    & ---   & ---  & ---
  & Seizure \& active task correctly lowest \\
\bottomrule
\end{tabular}
\end{table}

% ============================================================
\begin{thebibliography}{99}
% ============================================================

\bibitem{adams2019}
Adams, F.~C. (2019).
The degree of fine-tuning in our universe --- and others.
\textit{Physics Reports}, 807, 1--111.

\bibitem{bak1987}
Bak, P., Tang, C., \& Wiesenfeld, K. (1987).
Self-organized criticality.
\textit{Physical Review Letters}, 59(4), 381.

\bibitem{barabasi1999}
Barab{\'a}si, A.-L., \& Albert, R. (1999).
Emergence of scaling in random networks.
\textit{Science}, 286(5439), 509--512.

\bibitem{chaisson2001}
Chaisson, E.~J. (2001).
\textit{Cosmic Evolution: The Rise of Complexity in Nature}.
Harvard University Press.

\bibitem{crutchfield2001}
Crutchfield, J.~P., \& Shalizi, C.~R. (2001).
Thermodynamic depth of causal states: Objective complexity via minimal
representations.
\textit{Physical Review E}, 63(5), 054119.

\bibitem{england2013}
England, J.~L. (2013).
Statistical physics of self-replication.
\textit{The Journal of Chemical Physics}, 139(12), 121923.

\bibitem{gell-mann1996}
Gell-Mann, M., \& Lloyd, S. (1996).
Information measures, effective complexity, and total information.
\textit{Complexity}, 2(1), 44--52.

\bibitem{kauffman1993}
Kauffman, S.~A. (1993).
\textit{The Origins of Order: Self-Organization and Selection in Evolution}.
Oxford University Press.

\bibitem{kolmogorov1965}
Kolmogorov, A.~N. (1965).
Three approaches to the quantitative definition of information.
\textit{Problems of Information Transmission}, 1(1), 1--7.

\bibitem{ladyman2013}
Ladyman, J., Lambert, J., \& Wiesner, K. (2013).
What is a complex system?
\textit{European Journal for Philosophy of Science}, 3(1), 33--67.

\bibitem{langton1990}
Langton, C.~G. (1990).
Computation at the edge of chaos.
\textit{Physica D}, 42(1--3), 12--37.

\bibitem{lloyd1988}
Lloyd, S., \& Pagels, H. (1988).
Complexity as thermodynamic depth.
\textit{Annals of Physics}, 188(1), 186--213.

\bibitem{naszvadi2020}
Naszvadi, P. (2020).
Seeds is Turing-complete.
\textit{LifeWiki} / personal communication.
\url{https://conwaylife.com/wiki/Seeds}

\bibitem{prigogine1977}
Prigogine, I., \& Stengers, I. (1977).
\textit{Order Out of Chaos: Man's New Dialogue with Nature}.
Bantam Books.

\bibitem{tononi2004}
Tononi, G. (2004).
An information integration theory of consciousness.
\textit{BMC Neuroscience}, 5(1), 42.

\bibitem{watts1998}
Watts, D.~J., \& Strogatz, S.~H. (1998).
Collective dynamics of `small-world' networks.
\textit{Nature}, 393(6684), 440--442.

\bibitem{wolfram2002}
Wolfram, S. (2002).
\textit{A New Kind of Science}.
Wolfram Media.

\end{thebibliography}

\end{document}
